% SHARD-ID: AQM-16-PROJECTIVE-LINE-SHEAF-FIELD-CALCULATION
% SHARD-TITLE: Projective Line Sheaf Field Calculation
% SHARD-SUMMARY: Computes the finite sheaf-field construction for \(\mathbf P^1_{\mathbb F_3}\) and \(V=\mathbb F_3^2\).
% SHARD-SUMMARY: Lists the rational points, the \(H^0(\mathcal O(d))\) bases, point evaluations, scalar Gram forms, radicals, and Weyl counts.
% SHARD-SUMMARY: Records the exact Julia validator and its CSV outputs.
% SHARD-KEYWORDS: projective line, sheaf, O(d), finite field, symplectic target, Weyl labels

\section{Projective Line Sheaf Field Calculation}
\label{sec:projective-line-sheaf-field-calculation}

\subsection{Status and Local Evidence}

This shard is self-contained for the example
\[
  X=\mathbf P^1_{\mathbb F_3},\qquad V=\mathbb F_3^2.
\]
The algebraic geometry inputs are the local Stacks TeX files cited in
Shard \(AQM\)-15.  The exact finite calculation is checked by
\path{scripts/bridges/projective_line_sheaf_fields.jl}.  The AQM-16 output
files are
\path{runs/2026-05-24-projective-line-sheaf-fields/data/projective_line_sheaf_field_summary.csv}
and
\path{runs/2026-05-24-projective-line-sheaf-fields/data/projective_line_sheaf_field_basis_rows.csv}.

\subsection{The Field and the Points}

The field is
\[
  \mathbb F_3=\{0,1,2\},\qquad 2=-1,\qquad 1+1+1=0.
\]
The projective line has rational points given by nonzero pairs
\((a,b)\in\mathbb F_3^2\setminus\{(0,0)\}\) modulo nonzero scalar
multiplication:
\[
  [a:b]=[\lambda a:\lambda b]\quad(\lambda\in\mathbb F_3^\times).
\]
There are four classes.  The chosen representatives and order are
\[
  P_\infty=[1:0],\quad P_0=[0:1],\quad P_1=[1:1],\quad P_2=[2:1].
\]

\subsection{The Sheaves and Their Sections}

For every \(d\geq0\), the source computation gives
\[
  H^0(\mathbf P^1_{\mathbb F_3},\mathcal O(d))
  =
  \mathbb F_3[X,Y]_d.
\]
Thus a basis is
\[
  X^d,\ X^{d-1}Y,\ X^{d-2}Y^2,\ \ldots,\ Y^d.
\]
This is a statement about sections of \(\mathcal O(d)\), not about global
regular functions when \(d>0\).

For this finite rational-point calculation, evaluate \(X^{d-i}Y^i\) on the
fixed representatives above:
\[
  X^{d-i}Y^i(P_\infty)=
  \begin{cases}
  1,&i=0,\\
  0,&i>0,
  \end{cases}
\]
and
\[
  X^{d-i}Y^i(P_t)=t^{d-i}\qquad(t=0,1,2).
\]
This is the explicit chart-coordinate evaluation convention from
CONVENTIONS (r).  It is not silently canonical for every line bundle.

\subsection{Scalar Pairing and Symplectic Extension}

For scalar sections \(s,t\), define
\[
  B(s,t)=\sum_{P\in\mathbf P^1(\mathbb F_3)}s(P)t(P).
\]
All weights are \(1\).  The target is \(V=\mathbb F_3^2\) with
\(\omega((q,p),(q',p'))=pq'-qp'\).
The field space is
\[
  E_d=H^0(\mathbf P^1,\mathcal O(d))\otimes_{\mathbb F_3}V.
\]
If \(U_d=H^0(\mathbf P^1,\mathcal O(d))\), write
\[
  E_d=U_d q\oplus U_d p.
\]
For \((q,p),(q',p')\in U_d q\oplus U_d p\),
\[
  \Omega_d((q,p),(q',p'))=B(p,q')-B(q,p').
\]
If the scalar Gram matrix \(G_d=(B(e_i,e_j))\) has rank \(r\), then
\[
  \dim E_d=2(d+1),\quad
  \operatorname{rank}\Omega_d=2r,\quad
  \dim\operatorname{rad}(E_d)=2(d+1-r).
\]
The Weyl label space is \(E_{d,\rm red}=E_d/\operatorname{rad}(E_d)\).

\subsection{Degree 0: Constants}

The basis is \(1\).  Its value row is
\[
  1:\quad (1,1,1,1).
\]
Therefore
\[
  B(1,1)=1+1+1+1=4=1\in\mathbb F_3.
\]
So
\[
  G_0=[1].
\]
The scalar rank is \(1\), hence \(\Omega_0\) has rank \(2\).  The reduced
field has dimension \(2\), the Weyl Hilbert dimension is \(3\), and the Weyl
operator basis has \(9\) elements.

\subsection{Degree 1: Linear Sections}

The basis is \(X,Y\).  The value rows are
\[
  X:\ (1,0,1,2),\qquad
  Y:\ (0,1,1,1).
\]
Compute every scalar product:
\[
\begin{aligned}
B(X,X)&=1^2+0^2+1^2+2^2=1+0+1+4=6=0,\\
B(X,Y)&=1\cdot0+0\cdot1+1\cdot1+2\cdot1=3=0,\\
B(Y,Y)&=0^2+1^2+1^2+1^2=3=0.
\end{aligned}
\]
Thus
\[
  G_1=
  \begin{pmatrix}
  0&0\\
  0&0
  \end{pmatrix}.
\]
The two sections evaluate independently at the four points, but the chosen
all-weights-one scalar pairing is totally zero on their span.  Consequently
\[
  \operatorname{rank}\Omega_1=0,\qquad
  \operatorname{rad}(E_1)=E_1.
\]
The reduced Weyl label space is zero-dimensional.  This is the first concrete
warning that evaluation rank and symplectic rank are different questions.

\subsection{Degree 2: Quadratic Sections}

The basis is \(X^2,XY,Y^2\).  The value rows are
\[
\begin{array}{c|cccc}
&P_\infty&P_0&P_1&P_2\\
\hline
X^2&1&0&1&1\\
XY&0&0&1&2\\
Y^2&0&1&1&1
\end{array}
\]
The scalar products give
\[
  G_2=
  \begin{pmatrix}
  0&0&2\\
  0&2&0\\
  2&0&0
  \end{pmatrix}.
\]
For example,
\[
  B(X^2,Y^2)=1\cdot0+0\cdot1+1\cdot1+1\cdot1=2,
\]
and
\[
  B(XY,XY)=0+0+1+4=5=2.
\]
The determinant is
\[
  \det(G_2)=1\in\mathbb F_3,
\]
so \(G_2\) has rank \(3\).  Hence \(\Omega_2\) has rank \(6\), no radical,
and
\[
  \dim E_2=6,\qquad
  \dim\mathcal H=3^3=27,\qquad
  |\{W(e)\}|=3^6=729.
\]

\subsection{Degree 3: Cubic Sections}

The basis rows and Gram matrix are
\[
\begin{array}{c|cccc}
&P_\infty&P_0&P_1&P_2\\
\hline
X^3&1&0&1&2\\
X^2Y&0&0&1&1\\
XY^2&0&0&1&2\\
Y^3&0&1&1&1
\end{array}
\]
and
\[
  G_3=
  \begin{pmatrix}
  0&0&2&0\\
  0&2&0&2\\
  2&0&2&0\\
  0&2&0&0
  \end{pmatrix}.
\]
The validator computes rank \(4\), equivalently \(\det(G_3)=1\).  Thus
\[
  \dim E_3=8,\quad \operatorname{rank}\Omega_3=8,\quad
  \dim\mathcal H=3^4=81.
\]

\subsection{Degree 4: First Evaluation Kernel}

The scalar section dimension is \(5\), but there are only four rational
points.  The value rows include
\[
  X^3Y:\ (0,0,1,2),\qquad
  XY^3:\ (0,0,1,2),
\]
so
\[
  X^3Y-XY^3
\]
evaluates to zero at all four chosen representatives.  Therefore the scalar
evaluation map has a one-dimensional kernel.  The run confirms
\[
  \operatorname{rank}(\operatorname{ev})=4,\qquad
  \operatorname{rank}(G_4)=4.
\]
The scalar radical is exactly one-dimensional, and after tensoring with
\(V=\mathbb F_3^2\),
\[
  \dim\operatorname{rad}(E_4)=2.
\]
Thus
\[
  \dim E_4=10,\qquad
  \dim E_{4,\rm red}=8,\qquad
  \dim\mathcal H=3^4=81.
\]

\subsection{Run Summary}

\[
\begin{array}{c|c|c|c|c|c|c}
d&\dim U_d&\operatorname{rank}(\operatorname{ev})&
\operatorname{rank}G_d&\dim\operatorname{rad}E_d&
\dim E_{d,\rm red}&\dim\mathcal H\\
\hline
0&1&1&1&0&2&3\\
1&2&2&0&4&0&1\\
2&3&3&3&0&6&27\\
3&4&4&4&0&8&81\\
4&5&4&4&2&8&81
\end{array}
\]

The lessons are concrete: the projective sheaf field is controlled by
\(H^0(\mathcal O(d))\), the finite rational-point pairing is an additional
choice and can introduce radicals, and the Weyl-Heisenberg observable algebra
is formed after quotienting by this radical.
